% FIGURE LIST

%%%%%%% MAIN FIGURES %%%%%%%

\newcommand{\FigParasites}{
    \begin{wrapfigure}{R}{0.5\textwidth}
       \vspace{-16pt}
        \ffigbox[\FBwidth]
                {
                    \includegraphics[width=0.93\textwidth, keepaspectratio=true]{sheep-arc-figures/Fig 1 parasite counts.pdf}
                }
                {
                    \caption{\textbf{Vaccine- and age-mediated control of \emph{T. circumcincta} infection.} \textbf{a}, Worm burdens and \textbf{b}, total egg output in 3 month-old lambs with vaccination (3mo-Vax), 3 month-old lambs with adjuvant only (3mo-Ctrl), 6 month-old lambs with vaccination (6mo-Vax), and 6 month-old lambs with adjuvant only (6mo-Ctrl). Immunisation with the prototype vaccine led to a median 73.7\% reduction (P\textsubscript{vacc}~$=$~0.002, GLM) in worm counts and 50\% reduction (P\textsubscript{vacc}~$=$~0.026, GLM) in cumulative egg output.}
                     \label{fig:parasites}
                 }
       \vspace{-12pt}
    \end{wrapfigure}
}

% Alternative Fig 1
% \newcommand{\FigParasites}{
%     \begin{figure}[hp]

%         \includegraphics[width=0.6\textwidth, keepaspectratio=true]{sheep-arc-figures/Fig 1 parasite counts.pdf}

%         \caption{\textbf{Vaccine- and age-mediated control of \emph{T. circumcincta} infection.} \textbf{a}, Worm burdens and \textbf{b}, total egg output in 3 month-old lambs with vaccination (3mo-Vax), 3 month-old lambs with adjuvant only (3mo-Ctrl), 6 month-old lambs with vaccination (6mo-Vax), and 6 month-old lambs with adjuvant only (6mo-Ctrl). Immunisation with the prototype vaccine led to a median 73.7\% reduction (P\textsubscript{vacc}~$=$~0.002, GLM) in worm counts and 50\% reduction (P\textsubscript{vacc}~$=$~0.026, GLM) in cumulative egg output.}

%         \label{fig:parasites}

%     \end{figure}
% }

\newcommand{\FigCircle}{
\begin{figure}[hp]
    \includegraphics[width=\textwidth, keepaspectratio=true]{sheep-arc-figures/Fig 2 pathway circle.pdf}

    \caption{\textbf{Temporal dynamics of the expression of immune
        pathways that predict either post-mortem worm burden or
        cFEC identified by supervised machine learning
        ElasticNet.} Heat map indicates the pathway enrichment
        reported as negative $log_{10}^{\mathrm{P-value}}$,
        where the dark red denotes stronger enrichment within
        each pathway. Each concentric circle represents one of
        six time-points at which abomasal biopsies were taken
        after challenge infection. Worm burdens predicted by these
        pathways were measured at D49 post challenge.}

    \label{fig:circular}
\end{figure}

}

\newcommand{\FigEnetcorr}{
    \begin{figure}[htbp]
    \includegraphics[width=\textwidth, keepaspectratio=true]{sheep-arc-figures/Fig 3 Enet_corr.pdf}

    \caption{\textbf{Heat map of correlation coefficients} between
        \textbf{a}, worm burdens or \textbf{b}, cFEC and gene
        expression in the selected pathways before challenge
        (D0), seven (D7) and 21 (D21) days post challenge in
        vaccinated and control lambs. Colours represent negative
        (red) and positive (blue) Pearson correlation
        coefficients for each comparison. The immune pathways
        were clustered using k-means according to their
        correlation patterns over time in the vaccinated group.}

    \label{fig:enet_corr}
\end{figure}

}

\newcommand{\FigWGCNA}{
    \begin{figure}[htbp]
    \includegraphics[width=\textwidth, keepaspectratio=true]{sheep-arc-figures/Fig 4 WGCNA_gene_expression.pdf}

    \caption{\textbf{Gene expression within immune pathways
    significantly affected by immunisation, age, or both.}
        \textbf{a} Heat map of gene expression over time within
        the four treatments (3mo-Vax, 3mo-Ctrl, 6mo-Vax, and
        6mo-Ctrl). High to low scaled expression is denoted with
        blue to white hues. \textbf{b} Heat map of a generalised
        mixed model t value statistic (number of standard
        deviations from the mean)~\supercite{Bates2015} indicating
        the effect of age, vaccination, and the interaction
        between age and vaccination on the expression of each
        pathway. Purple indicates a dampening effect on the
        expression of the pathway while yellow indicates
        increased expression of that pathway relative to the
        global average. Pathways included are those presented in
        Fig.~\ref{fig:circular}. The full list of gene expression
        is available in Fig.~\ref{suppfig:PathwayDynamics}.}

    \label{fig:WGCNA_gene_expression}
\end{figure}
}

%%%%%%% SUPPLEMENTARY FIGURES %%%%%%%

\newcommand{\SuppFigTimeline}{
    \begin{figure}[H]
        \includegraphics[width=\textwidth, keepaspectratio=true]{sheep-arc-figures/Fig S1 exp timeline.pdf}

       \caption{\textbf{Experimental design.} Timeline of immunisation, challenge, and sampling of sheep.}

        \label{suppfig:timeline}
    \end{figure}
}

\newcommand{\SuppFigFlow}{
    \begin{figure}[H]
        \includegraphics[width=\textwidth, keepaspectratio=true]{sheep-arc-figures/Fig S2 analysis pipeline.pdf}

        \caption{\textbf{Data analysis pipeline overview.} Depiction of the
            flow through which raw transcript counts were taken
            to extract immunological information relevant to the
            response to vaccination and infection.
        }

        \label{suppfig:flow}
    \end{figure}
}

\newcommand{\SuppFigTSNE}{
    \begin{figure}[H]
    \includegraphics[width=\textwidth, keepaspectratio=true]{sheep-arc-figures/Fig S3 t_sne.pdf}

    \caption{\textbf{Unsupervised clustering of transcriptomes per
        treatment and day post challenge.} The $\sim$15K
        transcripts were reduced to two components using
        t-Distributed Stochastic Neighbour Embedding
        (t-SNE)~\supercite{Maaten2008} to visualise the datasets in
        two-dimensional space. Samples clustered weakly by
        treatment but more distinctly by day post challenge
        (DPC), specifically before \emph{vs.} after 14 DPC.}

    \label{suppfig:t-sne}
\end{figure}
}

\newcommand{\SuppFigEnetmodules}{
\begin{figure}[H]
    \includegraphics[width=\textwidth, keepaspectratio=true]{sheep-arc-figures/Fig S4 ElasticNet coefficients.pdf}

    \caption{\textbf{ElasticNet coefficients bar plot of WGCNA modules
        associated with worm burden and cFEC across 42 days post
        challenge.} Only modules with ElasticNet coefficients
        ≠ 0 are depicted and were retained for further
        analysis.
    }

    \label{suppfig:enetmodules}
\end{figure}
}

\newcommand{\SuppFigDiffexp}{
    \begin{figure}[H]
        \includegraphics[width=\textwidth, keepaspectratio=true]{sheep-arc-figures/Fig S5 diff_exp.pdf}

        \caption{\textbf{Time dynamics of differentially expressed (DE)
            genes.} \textbf{a}, Number of DE genes in each WGCNA
            modules within each of 4 pairwise DE pairwise
            comparisons: 3 month-old control \emph{vs.}
            6 month-old control (`3mo-Ctrl vs 6mo-Ctrl'),
            3 month-old vaccinated \emph{vs.} 3 month-old
              control (`3mo-Vax \emph{vs} 3mo-Ctrl'),
              6 month-old vaccinated \emph{vs.} 6 month-old
                control (`6mo-Vax \emph{vs} 6mo-Ctrl') and
              3 month-old vaccinated \emph{vs.} 6 month-old
                vaccinated (`3mo-Vax \emph{vs} 6mo-Vax').
                \textbf{b}, Venn diagram of all genes from four
                pairwise comparisons. \textbf{c} Heat map of the
                temporal expression patterns of significant DE
                gene expression. Individual genes within
                clusters 1—7 are shown in columns (see full
            gene list in Supplementary Data 2). Each row
        represents either 3mo or 6mo vaccinated lambs at
    different time points.}

        \label{suppfig:diff_exp}
    \end{figure}
}

\newcommand{\SuppFigPathwayDynamics}{
\begin{figure}[H]
    \includegraphics[width=\textwidth, keepaspectratio=true]{sheep-arc-figures/Fig S6 pathway_dynamics.pdf}

    \caption{\textbf{Time course of pathways significantly represented
        in abomasal transcriptomes.} Mean scaled expression
        levels in each pathway predictive of worm burdens and/or
    cFEC, selected in ~\nameref{fig:circular}. Points represent individual lambs at each time point and lines represent corresponding mean values for: vaccinated (solid lines); control (dashed lines); 3-month-old (red); and 6-month-old (blue) lambs.}

    \label{suppfig:PathwayDynamics}
\end{figure}
}
